\SourceLeaf{275}%
selbstständige Geometrie, die er gelegentlich einmal verfolgt und mit dem Namen Antieuklidische Geometrie bezeichnet habe.

\textsc{Gauss}, nach seiner öfters ausgesprochenen innersten Ansicht, betrachtete die drei Dimensionen des Raumes als eine specifische Eigenthümlichkeit der menschlichen Seele; Leute, welche diess nicht einsehen könnten, bezeichnete er einmal in seiner humoristischen Laune mit dem Namen Böotier. Wir können uns, sagte er, etwa in Wesen hineindenken, die sich nur zweier Dimensionen bewusst sind; höher über uns stehende würden vielleicht in ähnlicher Weise auf uns herabblicken, und er habe, fuhr er scherzend fort, gewisse Probleme hier zur Seite gelegt, die er in einem höhern Zustande später geometrisch zu behandeln gedächte.\}

\Needspace{6\baselineskip}
\begin{center}\textsc{Bemerkung.}\end{center}

Das von \textsc{Sartorius von Waltershausen} erwähnte Papier hat sich in \textsc{Gauss}' Nachlass vorgefunden, es bezieht sich jedoch im Wesentlichen nur auf die Erklärung des Zahlbegriffs und der vier Species.

\hfill \textsc{Stäckel}.
